% ============================================================
%  Chapter 1: Introduction
% ============================================================
%
% AIM: Motivate the project, state the problem, and outline the report.
%      Readable by someone with a general ML background but no CT knowledge.
%      ~2-3 pages.
%
% CONTENTS:
%   1.1 Clinical motivation for low-dose CT
%       - Why CT matters (medical imaging workhorse)
%       - Why dose reduction matters (radiation risk)
%       - The reconstruction quality vs dose trade-off
%
%   1.2 The reconstruction problem in one paragraph
%       - Forward model y = Ax + noise (informal, no equations yet)
%       - Classical methods struggle with sparse data
%       - Learned methods can help
%
%   1.3 This project
%       - Paper being reproduced: TFPnP (Wei et al., JMLR 2022)
%       - Why this paper: bridges PnP and RL, generic framework
%       - Implementation target: integrate into LION toolbox
%       - Dataset: specify which one you use
%
%   1.4 Contributions
%       - Reproduction of TFPnP sparse-view CT experiments
%       - Implementation within LION/tomosipo (different CT backend)
%       - Comparison against classical baselines (FBP, TV) and LPD
%       - [Optional] Extension to cone-beam geometry
%
%   1.5 Report outline
%       - One sentence per chapter
%
% FIGURES: None needed. Keep it clean prose.
%
% TIPS:
%   - Write this LAST (after all other chapters)
%   - Keep it short and punchy
%   - Don't repeat the abstract
% ============================================================

\chapter{Introduction}
\label{chap:introduction}

% TODO: Write after all other chapters are complete.